% ================================================
% 文件: 01_电子元器件.tex
% 章节: 第2章 电子元器件
% 所属部分: 电子元器件与材料
% 版本: v1.0
% ================================================
% 本章目录结构（树状）:
% \chapter{电子元器件}
% ├── \section{无源元件（电阻、电容、电感）}
% ├── \section{半导体器件（二极管、三极管）}
% ├── \section{电子管}
% ├── \section{集成电路}
% └── \section{传感器与换能器}
% ================================================

\chapter{电子元器件}
\label{chap:electronic_components}

\section{无源元件（电阻、电容、电感）}
\label{sec:passive_components}

无源元件是指不需要外部电源就能工作的电子元件。

\begin{itemize}
    \item \textbf{电阻器}：阻碍电流流动，单位为欧姆（$\Omega$）。电阻值由材料、长度和截面积决定：$R = \rho\frac{l}{S}$
    
    \item \textbf{电容器}：存储电荷，单位为法拉（F）。电容值：$C = \frac{\varepsilon S}{d}$
    
    \item \textbf{电感器}：阻碍电流变化，单位为亨利（H）。电感值与线圈匝数、磁芯材料有关
\end{itemize}

\section{半导体器件（二极管、三极管）}
\label{sec:semiconductor_devices}

半导体器件是以半导体材料为基础制成的电子元件。

\begin{itemize}
    \item \textbf{二极管}：具有单向导电性，由PN结构成。主要用于整流、检波等
    
    \item \textbf{三极管（BJT）}：具有电流放大作用，有NPN和PNP两种类型。共射极放大倍数$\beta = \frac{I_c}{I_b}$
    
    \item \textbf{场效应管（FET）}：电压控制器件，输入阻抗高。分为JFET和MOSFET
\end{itemize}

\section{电子管}
\label{sec:vacuum_tubes}

电子管是早期的电子放大器件，在真空环境中工作。

\begin{itemize}
    \item \textbf{二极管}：两个电极，用于整流
    \item \textbf{三极管}：三个电极（阴极、栅极、阳极），用于放大
    \item \textbf{四极管、五极管}：增加抑制栅极等，改善性能
\end{itemize}

电子管具有线性好、音质优美等优点，但体积大、功耗高。

\section{集成电路}
\label{sec:integrated_circuits}

集成电路（IC）是将大量电子元件集成在一小块半导体芯片上。

\begin{itemize}
    \item \textbf{SSI}：小规模集成电路（<100个元件）
    \item \textbf{MSI}：中规模集成电路（100-1000个元件）
    \item \textbf{LSI}：大规模集成电路（1000-100000个元件）
    \item \textbf{VLSI}：超大规模集成电路（>100000个元件）
\end{itemize}

\section{传感器与换能器}
\label{sec:sensors_transducers}

传感器是将非电信号转换为电信号的装置。

\begin{itemize}
    \item \textbf{温度传感器}：如热敏电阻、热电偶
    \item \textbf{光传感器}：如光敏电阻、光电二极管
    \item \textbf{压力传感器}：如压电传感器
    \item \textbf{换能器}：如扬声器、麦克风
\end{itemize}